\chapter{Putting Part III together}\label{p3-ch:revision}

\section{One chain of ideas}\label{p3-sec:chain}

The three sections of this part are one argument, and they close the argument begun in Parts~I and II:
\begin{enumerate}[itemsep=4pt]
\item Air holds water vapour up to a limit set by its temperature. Relative humidity compares the vapour with that limit; the mixing ratio and the dew point measure the vapour itself, and cooling air to its dew point makes dew, fog or cloud (Moisture~Q1--Q6, Chapters~\ref{p3-ch:moisture} and \ref{p3-ch:dew}).
\item Warm air columns are taller than cold ones (the scale-height law of Part~I, Chapter~\xref{p1-ch:pressure}), so pressure surfaces slope down from the warm tropics towards the cold poles, more steeply with height (Motion~Q1--Q6, Chapter~\ref{p3-ch:surfaces}). The uneven warming that causes this is the sun-angle effect of Part~II (Chapter~\xref{p2-ch:sun}).
\item The slope produces a pressure gradient force; on a rotating Earth the Coriolis force turns the moving air until the two balance and the wind blows along the isobars, the geostrophic wind (Motion~Q7, Chapter~\ref{p3-ch:geo}).
\item Near the ground friction slows the wind and turns it across the isobars towards low pressure, so air spirals into lows, rises, and forms cloud and rain (Motion~Q8--Q11, Chapter~\ref{p3-ch:fric}).
\item A weather map shows all of this at once: station models record the observations, streamlines show the flow, and isobars, fronts, cloud and rain show the weather systems it organises (Map~Q12--Q14, Chapters~\ref{p3-ch:station} to \ref{p3-ch:fronts}).
\end{enumerate}

\section{All the answers in one table}\label{p3-sec:summary}

\begin{table}[ht]
\centering
\caption{Every answer of Part~III, with the chapter where it is worked out.}\label{p3-tab:summary}
\small
\begin{tabularx}{\textwidth}{@{}lXr@{}}
\toprule
\textbf{Question} & \textbf{Answer} & \textbf{Ch.}\\
\midrule
Moisture Q1 & RH $=49.3$, $88.8$, $26.0$, $10.4$, $19.9\,\%$ & \ref{p3-ch:moisture}\\
Moisture Q2 & Humidity falls sharply: the vapour content is unchanged but the capacity of the warmer air is much larger & \ref{p3-ch:moisture}\\
Moisture Q3 & Mixing ratios $9.126$, $8.974$, $7.684$, $11.078$, $10.54\,\gkg$; ranks $3,4,5,1,2$ & \ref{p3-ch:moisture}\\
Moisture Q4 & RH does not say how much vapour is actually in the air & \ref{p3-ch:moisture}\\
Moisture Q5 & Dew point $\approx14\,\degC$ ($13.8$) & \ref{p3-ch:dew}\\
Moisture Q6 & $w_s=35.134\,\gkg$; $w=21.78\,\gkg$; dew point $\approx26\,\degC$ & \ref{p3-ch:dew}\\
Motion Q1 & Lines 5400$\to$5700\,m (500 mb) and 8700$\to$9500\,m (300 mb) & \ref{p3-ch:surfaces}\\
Motion Q2 & Saskatoon & \ref{p3-ch:surfaces}\\
Motion Q3 & Saskatoon is colder at each height; the cold column is thinner & \ref{p3-ch:surfaces}\\
Motion Q4 & z has the higher pressure: H at z (Monterrey side), L at x; arrow from z to x & \ref{p3-ch:surfaces}\\
Motion Q5 & At 9000\,m (800\,m against 140\,m of slope) & \ref{p3-ch:surfaces}\\
Motion Q6 & Wind speeds are greater at 9000\,m & \ref{p3-ch:surfaces}\\
Motion Q7 & North: H at 5520\,m, L at 5460\,m; south: H at 5640\,m, L at 5580\,m; force arrow to L, Coriolis arrow to H, wind between them & \ref{p3-ch:geo}\\
Motion Q8 & The Coriolis force is weakened, because it depends on wind speed & \ref{p3-ch:fric}\\
Motion Q9 & Red up (to low pressure), black up-right, blue down-right (right of wind), green down-left (against wind) & \ref{p3-ch:fric}\\
Motion Q10 & NH low, NH high, SH low, SH high (examples a to d) & \ref{p3-ch:fric}\\
Motion Q11 & On each box: red along the line to or from the centre; blue at right angles to the wind (right in the north, left in the south); green against the wind & \ref{p3-ch:fric}\\
Map Q12 & Model 1: 1002.4, 37, \textbf{21}, 0\,\%. Model 2: 999.8, 63, 60, thunder with intermittent rain. Model 3: 1011.7, 43, 20, 0\,\%, 25\,kt from SW, $+1.7$, rising. Model 4: 999.7, 27, 20, 100\,\%, continuous snow, 25\,kt from NW, $-1.1$, falling & \ref{p3-ch:station}\\
Map Q13 & Five or six smooth streamlines swirling anticlockwise round a centre in Iowa and Illinois & \ref{p3-ch:stream}\\
Map Q14 & Isobars at 4\,mb; L in northern Illinois/Iowa; cold front to the south-south-west, warm front to the east; cloud above 75\,\% wrapped round the low; rain north and east of the low and along the cold front & \ref{p3-ch:fronts}\\
\bottomrule
\end{tabularx}
\end{table}

\section{Mistakes to avoid}\label{p3-sec:mistakes}

\begin{enumerate}[label=\textbf{\arabic*.},itemsep=3pt]
\item \textbf{Ranking by relative humidity} instead of the actual mixing ratio (Moisture Q3).
\item \textbf{Reading the dew point from the wrong row} of Table~\ref{p3-tab:table1}, or at the wrong temperature: always check that the air could hold \emph{at least} the sample's vapour at the temperature you give (Moisture Q6: 26, not 25).
\item \textbf{Putting H and L the wrong way round} on the pressure surfaces: the lower pressure surface (colder column) is the L (Motion Q4).
\item \textbf{Drawing the Coriolis force to the wrong side:} right in the north, left in the south, always at right angles to the wind.
\item \textbf{Drawing the pressure gradient force away from low pressure.} It points \emph{towards} low pressure.
\item \textbf{Forgetting that the staff of a wind barb points to where the wind comes from} (Map Q12 and Q13).
\item \textbf{Mis-reading the pressure code:} 998 is 999.8, not 1099.8 or 998.0 (Map Q12).
\item \textbf{Isobars that cross,} or that stop in the middle of the map (Map Q14).
\end{enumerate}

\section{Self-test}\label{p3-sec:selftest}

Try these first. The answers follow.

\begin{tryit}[Self-test: one line each]
\begin{enumerate}[label=\textbf{S\arabic*.},leftmargin=*,itemsep=3pt]
\item Air at 20\,\degC\ has a mixing ratio of 7.5\,\gkg. What is its relative humidity?
\item That air is cooled to 10\,\degC. What is its relative humidity now, and what is its dew point?
\item Which holds more water vapour: air at 35\,\degC\ with RH 30\,\%, or air at 25\,\degC\ with RH 60\,\%?
\item Find the dew point of air at 30\,\degC\ and RH 50\,\%, and the height at which cloud would form if it were lifted.
\item True or false: the dew point can be higher than the air temperature. Explain.
\item Why do pressure surfaces slope down towards cold air?
\item The 300\,mb surface is at 8700\,m over place S and 9500\,m over place M. In which direction does the pressure gradient force act at 9000\,m?
\item Two 500\,hPa height lines 90\,m apart are 600\,km apart at 35\deg N. Find the geostrophic wind.
\item In the Southern Hemisphere a wind blows towards the west. Where is the low pressure, to its north or south?
\item What does friction do to (a) the speed and (b) the direction of the surface wind?
\item Over land the surface wind is 8\,\ms\ and crosses the isobars at 40\deg. Estimate the geostrophic wind.
\item Air spirals inward and clockwise. Which hemisphere, which pressure system?
\item Decode the pressure codes 087 and 962.
\item Convert 50\,\degF\ and 14\,\degF\ to \degC.
\item A wind barb has two full barbs and one pennant. What is the speed?
\item Behind a front the temperature falls 18\,\degF\ and the wind turns clockwise from south-west to north-west. Which kind of front?
\item A staff points to the north-east. In which direction does the air move?
\item The 1012\,mb isobar lies between stations at 1009.6 and 1013.2\,hPa. Where?
\item Why does a strong wind go with closely spaced isobars?
\item What can relative humidity not tell you?
\end{enumerate}
\end{tryit}

\begin{answerbox}[Answers to the self-test]
\begin{enumerate}[label=\textbf{S\arabic*.},leftmargin=*,itemsep=3pt]
\item $7.5/14.956=50\,\%$.
\item $w_s(10)=7.762$: RH $=7.5/7.762=97\,\%$. Dew point: Table~\ref{p3-tab:table1} gives 6.771 at 8\,\degC\ and 7.762 at 10\,\degC, so $T_d=8+2(7.5-6.771)/(7.762-6.771)=9.5\,\degC$.
\item $w_s(35)\approx37.3$ gives $0.30\times37.3=11.2\,\gkg$; $w_s(25)\approx20.5$ gives $0.60\times20.5=12.3\,\gkg$. The 25\,\degC\ air holds more.
\item $w=0.5\times27.694=13.85\,\gkg$: between 18\,\degC\ (13.162) and 20\,\degC\ (14.956), $T_d=18+2(13.85-13.162)/1.794=18.8\,\degC$. Cloud base $\approx125\times(30-18.8)=1400$\,m.
\item False (in open air): it would mean the air is oversaturated and the surplus vapour would condense at once. At most $T_d=T$ at saturation.
\item Cold air is denser, so a column of it is thinner: the same pressure drop uses up less height, so the pressure surfaces are lower over cold air.
\item From the high (the 300\,mb surface is higher over M, so z over M has the higher pressure) to the low: from M towards S.
\item $f=2\times7.292\times10^{-5}\sin35\deg=8.37\times10^{-5}$; $V_g=\dfrac{9.81\times90}{8.37\times10^{-5}\times6\times10^{5}}=17.6\,\ms$ (34 kt).
\item To the \textbf{north}: in the Southern Hemisphere the low is on the right of the wind (Buys-Ballot), and the right of a westward wind is north.
\item (a) It reduces the speed. (b) It turns the wind across the isobars towards low pressure.
\item $V=V_g\cos\alpha$, so $V_g=8/\cos40\deg=10.4\,\ms$.
\item Inward and clockwise: Southern Hemisphere low.
\item $087\to1008.7$\,hPa; $962\to996.2$\,hPa.
\item $50\,\degF=10\,\degC$; $14\,\degF=-10\,\degC$.
\item $2\times10+50=70$ knots.
\item A cold front.
\item The air moves towards the south-west.
\item $(1012-1009.6)/(1013.2-1009.6)=0.67$: two thirds of the way from the lower to the higher station.
\item Close isobars mean a large pressure gradient force, which needs a large wind to balance it with the Coriolis force (and friction).
\item How much vapour is actually in the air (or how the air compares with air at another temperature): it only compares the vapour with the capacity at the present temperature.
\end{enumerate}
\end{answerbox}

\section{A checklist}\label{p3-sec:checklist}

\begin{itemize}[label=$\square$]
\item I can look up, and interpolate, in Table~\ref{p3-tab:table1}.
\item I can find an RH, a mixing ratio and a dew point from any two of $T$, $w$, RH.
\item I can explain why pressure surfaces slope and why wind increases with height.
\item I can draw the three forces correctly in both hemispheres.
\item I can name the four systems of Figure~\ref{p3-fig:fig6} from their rotation.
\item I can decode a station model: pressure, temperature, dew point, sky, wind, weather, tendency.
\item I can draw streamlines (arrows downwind) and isobars (interpolating, never crossing).
\item I can say how to recognise a front, and which way the wind turns across it.
\end{itemize}
